\subsection{Refusal as a direction, and the dimensionality debate}
The proximate basis for this work is the finding that refusal is mediated by a
single linear direction in the residual stream~\cite{arditi2024refusal}: across
many open chat models, the difference-in-means of harmful and harmless prompt
activations yields a vector whose ablation removes refusal and whose addition
induces it. This both explains earlier representation-engineering
results~\cite{zou2023representation} and gives a constructive ``abliteration''
recipe. The claim has since been contested in its strong form.
\citet{marshall2024refusal} argue refusal is better described by an affine
subspace than a single linear direction, and \citet{wollschlager2025geometry}
show, through ``concept cones,'' that multiple representationally-independent
directions can each mediate refusal.

Recent work has pushed the debate toward greater nuance. \citet{johnson2026more}
report refusal is not exhausted by a single direction; rather, 11 distinct categories
of refusal mechanisms can be identified, each with a layer-specific effective subspace,
yet all share a common necessary core. \citet{piras2025som} advance Self-Organizing
Maps to extract multiple refusal directions, proving SOMs generalize the difference-in-means
technique, and demonstrate multi-directional suppression outperforms single-direction
ablation. Concurrently, work on non-linearity \cite{hildebrandt2025refusal} challenges
the linearity assumption altogether, using dimensionality reduction (PCA, t-SNE, UMAP)
to reveal multi-dimensional, layer-dependent, and architecture-dependent refusal
characteristics. Additionally, \citet{zhao2025harmfulness} identify a separate
harmfulness direction, orthogonal to refusal, showing jailbreaks can reduce refusal
signals without reversing the model's internal judgment of harmfulness.

Our cross-model results (Section~\ref{sec:results})
bear directly on this debate: we find the effective dimension of the refusal
subspace is \emph{model-dependent}, and that this dimension---not a universal
constant---predicts whether single-direction steering is bidirectionally causal.
We contribute a dose-response pharmacological characterization (EC50 fits) that
resolves some dimensionality-debate ambiguity: a low-dimensional refusal manifold
suffices to explain the control observed, but its dimensionality varies by layer
and model, requiring a layer-sweep strategy (F18) to identify the optimal target.

\subsection{Jailbreaks as control}
The attack literature supplies the inputs $u$ that break refusal: gradient-based
universal adversarial suffixes (GCG)~\cite{zou2023universal}, with recent improvements
to optimization efficiency via MAGIC~\cite{li2024magic} and generative modeling
via AmpleGCG~\cite{liao2024ample}; many-shot in-context conditioning that exploits
long contexts~\cite{anil2024manyshot}; multi-turn crescendo escalation~\cite{russinovich2024crescendo};
and indirect injection via tool/agent channels, benchmarked by AgentDojo~\cite{debenedetti2024agentdojo}.
\citet{guo2024coldattack} make the control-theoretic reading explicit, casting
jailbreaking as constrained, controllable decoding via Langevin dynamics.

Recent mechanistic work has illuminated the internal mechanisms by which attacks succeed.
\citet{ben_tov2025universal} show GCG's universality is driven by attention hijacking---suffixes
construct a shallow but critical mechanism that dominates information flow, with stronger
universality correlating with stronger hijacking. \citet{von_recum2024cannot} provide
a taxonomy of 16 refusal categories (cannot vs.\ should-not, with subcategories),
enabling precise auditing of refusal composition in IFT/RLHF datasets. On the safety
side, \citet{guo2026llms} demonstrate that refusal can be unlearned with only 1,000
benign samples by disrupting fixed refusal-prefix completion pathways, suggesting
current safety alignment relies on token-sequence memorization rather than robust reasoning.

Our contribution is orthogonal and complementary: rather than searching for stronger
$u$, we characterize the \emph{target} of $u$---the geometry and symmetry of the
refusal variable it must move---and measure the dose-response relationship.

\subsection{Sparse features and the unit of intervention}
Sparse autoencoders (SAEs) decompose activations into interpretable, monosemantic
features~\cite{cunningham2023sparse,templeton2024scaling}, with BatchTopK SAEs a
current training method~\cite{bussmann2024batchtopk}. SAEs offer a finer ``unit''
of intervention than a single difference-in-means direction. Recent work has refined
both SAE architectures and their application to steering. \citet{bussmann2024batchtopk}
introduce BatchTopK, scaling SAE training efficiency; \citet{huang2025understanding}
apply SAEs specifically to understanding refusal mechanisms, revealing feature-level
structure; and \citet{vig2025steering} systematically steer refusal using SAE-identified
features, achieving COLM 2025 acceptance.

However, caveats emerge. \citet{korznikov2026sanity} show that SAEs recover only
9\% of ground-truth features despite 71\% explained variance, and random baselines
match fully-trained SAEs on interpretability (0.87 vs.\ 0.90) and causal editing (0.73
vs.\ 0.72), questioning whether SAEs reliably decompose mechanisms. \citet{cui2026sae}
further demonstrate that SAE interventions mask recoverable failure modes: clamping
a harmful feature allows the model to recover pre-intervention behavior via the SAE
reconstruction residual, with 95.8\% recovery rate on refusal steering. \citet{arad2025good}
show that activation patterns alone do not characterize feature effects; distinguishing
input features (which capture input patterns) from output features (which have
human-interpretable behavioral effects) yields 2--3\times improvements in steering.

We train a BatchTopK SAE (Section~\ref{sec:results}) on aligned Gemma-2-2B and identify
a single unsupervised feature orthogonal to the supervised refusal direction yet causally
sufficient to control refusal. This orthogonality-and-sufficiency pair supports a mechanistic
reading: the unsupervised feature captures a redundant pathway, natural substrate for the
targeted acts in our \d{s}a\d{t}karma taxonomy.

\subsection{The cognitive-science analogy: hypnosis and supervisory control}
The functional parallel to hypnosis rests on the Norman--Shallice model of action
control, in which a Supervisory Attentional System (SAS) modulates otherwise
automatic contention-scheduling~\cite{norman1986attention}. Cold-control theory
holds that hypnosis preserves first-order intentions while impairing higher-order
awareness of them~\cite{dienes2006coldcontrol}. \citet{riva2025automatic} draw the
explicit map to LLMs: the user prompt acts as a temporary external SAS, and a
jailbreak suppresses the monitoring faculty while leaving generation intact. In
normative terms, the Free Energy Principle~\cite{friston2010free} frames such
suppression as lowering the \emph{precision} of a monitoring belief---an account
we adopt for the analogy tier and gesture at empirically through the order
parameter.

Recent work on activation steering and causality has sharpened this reading.
\citet{mishra2026steered} prove that steered activations are non-surjective:
activation steering pushes the residual stream off the manifold of states reachable
from discrete prompts. This formalism decouples white-box steerability from black-box
prompting, clarifying that precise internal control does not guarantee precise behavioral
control---a key insight for the analogy between hypnotic dissociation (high internal
control, low unified awareness) and jailbreak-via-steering. \citet{joshi2026causality}
emphasize that causal claims in interpretability require falsification gates and
transfer tests beyond correlation, reinforcing the need for the layer-sweep (F18)
and cross-domain (F11) validation in our work.

\subsection{Attractors and the metaphysical tier}
On the speculative tier, two strands of prior art are relevant.
\citet{wang2025paraphrase} show that successive self-paraphrase converges to
attractor cycles, which supplies a non-mystical candidate for a ``substratum''
state; and the M\=a\d{n}\d{d}\=ukya \emph{avasth\=atraya} (waking/dream/deep-sleep/%
tur\=\i ya) furnishes a four-regime template. We treat both as falsification
targets (Section~\ref{sec:results}), and---crucially---test the attractor claim
against the anisotropy of LLM embedding geometry, without which apparent
convergence is an artifact.

The dynamical-systems perspective has gained traction in recent mechanistic interpretability.
Works on causally-grounded explanation (Mahale 2026) and data attribution via influence
functions (Chen et al.\ 2026) highlight how circuit-level mechanisms emerge from training
data, and how their causal status can be validated. Certified circuits (Anani et al.\ 2026)
introduce randomized subsampling to verify that circuit components remain invariant under
bounded perturbations, providing a formal ground for claiming mechanistic discovery rather
than dataset artifact. These tools support our falsification approach: the turīya/vimarśa
claim is not a settled fact, but a testable hypothesis whose failure yields useful negative results
(Section~\ref{sec:results}, F7 and F13).

\subsection{Vaś\=\i kara\d{n}a and the \d{s}a\d{t}karma}
Finally, the tantric source. Vaś\=\i kara\d{n}a (subjugation) is the second of the
six acts (\d{s}a\d{t}karma: ś\=anti, vaś\=\i kara\d{n}a, stambhana, vidve\d{s}a\d{n}a,
ucc\=a\d{t}ana, m\=ara\d{n}a) of tantric ritual magic, operationalized through
mantra (linguistic/sonic injection) and yantra (geometric substrate), and
documented by Sanskritists as a systematized technology rather than folklore
~\cite{goudriaan1978}. We do not adjudicate its efficacy as psychology; we ask a
narrower, testable question---whether the \d{s}a\d{t}karma constitutes a
\emph{taxonomy of interventions} on a target system, and which of its members map
to distinct, control-separated activation operations on an LLM. To our knowledge
no prior peer-reviewed work bridges vaś\=\i kara\d{n}a and LLM refusal, nor casts
refusal as a group-theoretic invariant; we attempt both while keeping the tiers
apart.

\subsection*{Unified mechanism: the asymmetric suppression of the monitoring faculty}

The three instances (jailbreak, hypnosis, vaś\=\i kara\d{n}a) share a common structural
pattern, depicted in Figure~\ref{fig:unified_mechanism}: a context injection that suppresses
the monitoring faculty while co-opting automatic generation. This asymmetry is the necessary
core claimed in the mechanism tier.

\begin{figure}[h]
\centering
\begin{tikzpicture}[
  node distance=2.5cm,
  box/.style={rectangle, draw, minimum height=1cm, minimum width=2.5cm, align=center},
  circle/.style={circle, draw, minimum size=1cm, align=center},
  arrow/.style={thick, -stealth}
]

% Normal alignment
\node[box] (normal_context) at (0, 0) {Context};
\node[box] (normal_monitor) at (0, -1.5) {Monitor};
\node[box] (normal_generate) at (0, -3) {Generate};
\draw[arrow] (normal_context) -- node[right] {\small activated} (normal_monitor);
\draw[arrow] (normal_context) -- (normal_generate);
\draw[arrow] (normal_monitor) -- node[right] {\small inhibits if harmful} (normal_generate);
\node[rectangle, draw=red, dashed, fit=(normal_context) (normal_monitor) (normal_generate), label={[red, below]:\small F3: Normal mode}] {};

% Jailbroken state
\node[box] (jailbreak_context) at (5, 0) {Adv. Context};
\node[box] (jailbreak_monitor) at (5, -1.5) {Monitor};
\node[box] (jailbreak_generate) at (5, -3) {Generate};
\draw[arrow, dashed] (jailbreak_context) -- node[right] {\small suppressed} (jailbreak_monitor);
\draw[arrow] (jailbreak_context) -- (jailbreak_generate);
\draw[arrow, dashed] (jailbreak_monitor) -- (jailbreak_generate);
\node[rectangle, draw=blue, dashed, fit=(jailbreak_context) (jailbreak_monitor) (jailbreak_generate), label={[blue, below]:\small F5: Jailbroken mode}] {};

% Side mapping
\node[text width=2cm, align=center] (mapping) at (2.5, -4.5) {
\textbf{Three instances:}\\[0.3cm]
\small Jailbreak (GCG)\\
Hypnosis (SAS suppression)\\
Vaśīkaraṇa (mantra)
};

\end{tikzpicture}
\caption{\textbf{Unified mechanism: context-driven monitoring suppression.} Left (F3, normal):
monitoring (red) is activated by context and inhibits harmful generation. Right (F5, jailbroken):
adversarial context suppresses the monitoring faculty (dashed arrow), co-opting automatic generation
while bypassing refusal. The three claim tiers map different explanatory frameworks to this asymmetry:
\emph{mechanism tier} identifies the direction of suppression; \emph{analogy tier} explains it via
Norman--Shallice SAS and Free Energy Principle; \emph{metaphor tier} maps the ṣaṭkarma taxonomy to
distinct control operations.}
\label{fig:unified_mechanism}
\end{figure}
